\SourceLeaf{222}%
\bigskip\Needspace{10\baselineskip}
\pdfbookmark[1]{Fortgesetzte Nachrichten. Pallas Olbersiana. 1802 Juli}{pallas180207}
\begin{center}
{\small FORTGESETZTE NACHRICHTEN ÜBER DEN NEUEN HAUPT-PLANETEN\par}
PALLAS OLBERSIANA.\par
\medskip\hrule\medskip
Monatliche Correspondenz. Band VI. S.\ 71—96. 1802 Juli.\par
\medskip\hrule
\end{center}\medskip
—\quad —\quad —

„Der sehr lehrreiche Aufsatz vom Prof.\ \textsc{Wurm} in Ihrem Junius-Hefte (schreibt Dr.\ \textsc{Gauss} unterm 26. Junius) zeigt leider die Ungewissheit, die noch bei der Masse des Jupiter Statt findet, dass wir die Störungs-Gleichungen der \textit{Ceres} und der \textit{Pallas} noch nicht mit der Zuverlässigkeit werden bestimmen können, die wol zu wünschen wäre. Dagegen werden die Beobachtungen der \textit{Ceres}, und noch mehr die der \textit{Pallas} auch den grossen Nutzen haben, dass sich daraus nach einem oder ein Paar Umläufen jene Masse sehr gut wird bestimmen lassen, welches dann auf viele andere wichtige Punkte der Astronomie nicht anders als wohlthätig zurückwirken kann. So bietet in dieser erhabenen Wissenschaft, die der Himmel dem Menschen so recht zur Cultur seiner Kräfte, und zur Erhebung über das Irdische geschenkt zu haben scheint, immer eine Entdeckung der andern die Hand; die grossen Entdeckungen stehen niemals allein, sondern sind oft eben so wichtig in ihren Folgen, als an sich selbst. Eben jene Unbestimmtheit wird mir noch eine neue Veranlassung sein, meinem vormaligen schon erwähnten Vorsatze getreu zu bleiben; nemlich ausser der, mit genauester Rücksicht auf die Störungen Jupiters wenigstens zu berechnenden neuen Ellipse für die \textit{Ceres} noch eine andere, auf eben dieselben Beobachtungen gegründete \textit{reine} Ellipse ohne alle Störungen zu berechnen. Ich glaube, es müsste doch interessant sein zu sehen, in wie fern sich künftiges Jahr schon eine \textit{gewisse} Spur von den Einwirkungen Jupiters in die Bewegung der \textit{Ceres} zeigen wird, denn bis jetzt zeigt sich noch nicht die allergeringste, sondern die Beobachtungen von anderthalb Jahren werden sich ohne Zweifel recht gut ohne Zwang durch eine \textit{reine Ellipse} darstellen lassen. Dass die VII. Elemente im Mai $40''$ abweichen, darüber darf man sich gar nicht wundern, wenn man bedenkt, auf wie dürftige Beobachtungen sie gegründet waren. Ich hatte ja nur die Palermer Beobachtungen und ein Paar von Ihnen, wozu die Abweichungen fehlten. Aus dieser Ursache kann ich auch meine bisherige Meinung noch nicht fahren lassen, dass es eine ganz unnöthige Vermehrung der Arbeit gewesen sein würde, wenn ich schon damals auf die Störungen hätte Rücksicht nehmen wollen. Ich sehe in der That nicht, dass dadurch bisher etwas besser hätte gemacht werden können; wohl aber glaube ich, dass vieles schlechter gemacht wäre. Ich wenigstens würde schwerlich Geduld genug gehabt haben, eine so grosse Anzahl Beobachtungen mit den Elementen haarscharf zu vergleichen und eine Ephemeride für eine geraume Zeit zu berechnen, wenn ich noch jedesmal eine so grosse Anzahl von Gleichungen hätte berechnen müssen. Eine ohne Rücksicht auf die Störungen berechnete Ellipse, die sich genau an die Beobachtungen hält, muss eine geraume Zeit diese Störungen selbst mit einschliessen, und diese Ellipse muss uns, so lange es nur darauf ankommt, theils die Beobachtungen zu erleichtern, theils die gemachten zu erörtern, und ihren respectiven Werth zu würdigen, weit wichtiger sein, als eine \textit{mittlere Ellipse}, die erst dann interessant sein kann, wenn die Dauer der Beobachtungen lang genug ist, um sie mit einer solchen Sicherheit bestimmen zu können, dass man etwa eine Anzahl Jahre damit rückwärts gehen kann, um eine Nachsuchung in den ältern Sternverzeichnissen anzustellen. Jetzt kann man sich nun schon nach anderthalb Jahren etwas ziemlich zuverlässiges davon versprechen, und ich werde daher künftig, wenn ich erst den Schluss der Palermer
\SourceLeaf{223}%
Beobachtungen habe, diese Rechnungen mit aller mir möglichen Sorgfalt anstellen. Aus diesen Gründen glaube ich auch, dass es wenig Nutzen haben würde, bei der \textit{Pallas} in diesem Jahre schon auf die Störungen Rücksicht zu nehmen, zumal da meiner Meinung nach die bei andern Planeten bisher üblichen Methoden bei diesem Planeten keinesweges die nöthige Schärfe geben würden. Nur dann, wenn die Lichtschwäche des Planeten, \textit{quod Deus avortat}, künftiges Jahr dessen Wiederauffindung vereiteln sollte, und man also zwei Jahre mit den Beobachtungen von diesem Jahre hinaus gehen müsste, würde es gut sein, von der Einwirkung des Jupiter Rechnung zu tragen, wiewol meines Erachtens auf einem andern als dem gewöhnlichen Wege.“\quad —\quad —\quad —

\Needspace{14\baselineskip}
\begin{center}
\textit{Elemente der Pallas Olbersiana aus 42tägigen Seeberger und Greenwicher Meridian-Beobachtungen.}\par\medskip
(III.)\par\medskip
{\small\setlength{\tabcolsep}{2pt}
\begin{tabular}{l r}
Epoche 1802 März 31 Mittag in Seeberg $\ .\ .\ .\ .\ .\ .\ .$&$162^0\ 25'\ 45''9$\\
tägliche mittlere tropische Bewegung $\ .\ .\ .\ .\ .\ .\ .$&$769''547$\\
Logarithmus der halben grossen Axe $\ .\ .\ .\ .\ .\ .\ .\ .$&$0.4425664$
\end{tabular}
\par\smallskip
$\left.\begin{array}{l}\text{Sonnenferne}\\\text{Knoten}\quad .\end{array}\right\}\ \text{für die Epoche und siderisch ruhend}\left\{\begin{array}{lr}
.\ .&300^0\ 58'\ 47''7\\
.\ .&172^0\ 28'\ 17''9
\end{array}\right.$
\par\smallskip
\begin{tabular}{l r}
Excentricität $\ .\ .\ .\ .\ .\ .\ .\ .\ .\ .\ .\ .\ .\ .\ .\ .\ .$&$0.2476402$\\
Neigung $\ .\ .\ .\ .\ .\ .\ .\ .\ .\ .\ .\ .\ .\ .\ .\ .\ .\ .$&$34^0\ 39'\ 10''7$
\end{tabular}}
\end{center}

Es ist doch höchst merkwürdig, dass die mittlere tropische Bewegung der \textit{Pallas} jener der \textit{Ceres}, unabsichtlich und so ganz von selbst, immer näher kommt. Nach den IV. Elementen der Ceres-Bahn (December-Heft 1801 S.\ 647 [203 d.\ B.]) bestimmte Dr.\ \textsc{Gauss} diese mittlere Bewegung auf $770''914$; diese ist von obiger bei der \textit{Pallas} nur anderthalb Secunden verschieden. Allein da auf beiden Bewegungen noch einige Ungewissheit haftet, so kann man jetzt eigentlich noch gar keinen Unterschied angeben.

Mit obigen III. Elementen verglich nun Dr.\ \textsc{Gauss} die Reihe nachstehender Beobachtungen, und fand ihre herrliche Übereinstimmung, wie folgende Tabelle zeigt.
\begin{center}
{\small\setlength{\tabcolsep}{4pt}
\begin{tabular}{r r|r|r|r|r|l}
\multicolumn{2}{c|}{1802}&\shortstack{Berechnete\\gerade Aufsteig.\\der Pallas}&\shortstack{Berechnete\\Abweichung\\der Pallas}&\multicolumn{2}{c|}{Unterschied}&Beobachter.\\
\multicolumn{2}{c|}{}&&&in AR.&in Declin.&\\\hline
April&4&$183^0\ 44'\quad5''6$&$13^0\ 54'\ 53''5$&$-1''0$&$+1''5$&von Zach\\
&5&$183\ 34\ 24.4$&$14\ 13\ 20.2$&$+0.$&$-2.7$&von Zach\\
&7&$183\ 15\ 38.5$&$14\ 48\ 59.3$&$0.7$&$-2.8$&von Zach\\
&8&$183\quad6\ 33.9$&$15\quad6\ 12.8$&$-3.9$&$+2.8$&von Zach\\
&10&$182\ 48\ 56.3$&$15\ 39\ 43.5$&$-1.7$&$+1.5$&Dr.\ Burckhardt\\
&12&$182\ 32\ 33.7$&$16\ 11\quad6.4$&$+2.2$&$+8.4$&Méchain\\
&13&$182\ 24\ 46.8$&$16\ 26\quad8.3$&$+2.4$&$+2.3$&Méchain\\
&15&$182\ 10\ 14.1$&$16\ 54\ 33.7$&$-2.4$&$+2.9$&von Zach\\
&15&$182\ 10\quad4.2$&$16\ 54\ 53.3$&$+1.3$&$+1.9$&Méchain\\
&16&$182\quad3\quad9.7$&$17\quad8\ 36.3$&$+1.7$&$-0.5$&Méchain\\
&17&$181\ 56\ 33.7$&$17\ 21\ 52.8$&$+3.4$&$+3.0$&Méchain\\
&18&$181\ 50\ 25.0$&$17\ 34\ 25.3$&$-5.6$&$\ .\ .\ .$&von Zach\\
&19&$181\ 44\ 26.4$&$17\ 46\ 50.0$&$+1.1$&$-4.4$&von Zach\\
&23&$181\ 23\ 41.9$&$18\ 32\ 27.4$&$-0.1$&$-0.6$&Dr.\ Maskelyne\\
&24&$181\ 19\ 32.4$&$18\ 42\ 25.2$&$-5.4$&$\ .\ .\ .$&von Zach\\
&25&$181\ 15\ 35.3$&$18\ 52\ 15.9$&$+3.1$&$\ .\ .\ .$&von Zach\\
&25&$181\ 15\ 28.6$&$18\ 52\ 33.1$&$+0.3$&$+3.4$&Dr.\ Maskelyne\\
&26&$181\ 11\ 59.4$&$19\quad1\ 41.6$&$+2.4$&$-7.8$&von Zach\\
&26&$181\ 11\ 53.3$&$19\quad1\ 58.1$&$+3.1$&$-4.5$&Dr.\ Maskelyne\\
&27&$181\quad8\ 44.8$&$19\ 10\ 42.7$&$+0.8$&$-4.1$&von Zach
\end{tabular}}
\end{center}
\expandafter\def\csname HistoricalNoteText20\endcsname{Bei der \textit{Greenwicher} Beobachtung der AR den 7. Mai scheint irgend ein Versehen vorgefallen zu sein, weil sie etwa $39''$ zu klein scheint.}
\SourceLeaf{224}%
\begin{center}
{\small\setlength{\tabcolsep}{4pt}
\begin{tabular}{r r|r|r|r|r|l}
\multicolumn{2}{c|}{1802}&\shortstack{Berechnete\\gerade Aufsteig.\\der Pallas}&\shortstack{Berechnete\\Abweichung\\der Pallas}&\multicolumn{2}{c|}{Unterschied}&Beobachter\\
\multicolumn{2}{c|}{}&&&in AR.&in Declin.&\\\hline
April&29&$181^0\quad3'\ 20''9$&$19^0\ 27'\ 31''5$&$+4''3$&$-12''2$&von Zach\\
&30&$181\quad1\ 12.1$&$19\ 35\ 19.6$&$+2.2$&$-13.9$&von Zach\\
Mai&1&$180\ 59\ 24.6$&$19\ 42\ 44.8$&$+2.1$&$+4.0$&von Zach\\
&1&$180\ 59\ 21.7$&$19\ 42\ 57.7$&$+2.8$&$+3.3$&Dr.\ Maskelyne\\
&2&$180\ 57\ 59.7$&$19\ 49\ 46.5$&$+7.8$&$-5.5$&von Zach\\
&2&$180\ 57\ 57.5$&$19\ 49\ 58.7$&$+0.2$&$+2.0$&Dr.\ Maskelyne\\
&3&$180\ 56\ 56.9$&$19\ 56\ 25.9$&$+1.3$&$-21.8$&von Zach\\
&4&$180\ 56\ 15.0$&$20\quad2\ 53.0$&$+2.1$&$-0.3$&Dr.\ Maskelyne\\
&5&$180\ 55\ 56.9$&$20\quad8\ 36.2$&$+2.3$&$-2.5$&von Zach\\
&6&$180\ 56\quad0.3$&$20\ 14\quad8.8$&$+6.2$&$\ .\ .\ .$&von Zach\\
&7&$180\ 56\ 24.7$&$20\ 19\ 20.4$&$+0.9$&$-14.9$&von Zach\\
&7&$180\ 56\ 25.8$&$20\ 19\ 29.4$&$+39.3$&$-0.5$&Dr.\ Maskelyne\\
&8&$180\ 57\ 13.0$&$20\ 24\ 10.8$&$+3.0$&$-9{,}3$&von Zach\\
&9&$180\ 58\ 23.9$&$20\ 28\ 49.3$&$+2.3$&$+6.8$&Dr.\ Maskelyne\\
&11&$181\quad1\ 44.6$&$20\ 36\ 41.9$&$+2.9$&$-31.3$&von Zach\\
&11&$181\quad1\ 48.3$&$20\ 36\ 48.5$&$+3.4$&$+0.3$&Dr.\ Maskelyne\\
&13&$181\quad6\ 37.6$&$20\ 43\ 31.4$&$+4.6$&$-1.7$&Dr.\ Maskelyne\\
&14&$181\quad9\ 33.4$&$20\ 46\ 25.2$&$+4.0$&$-1.5$&Dr.\ Maskelyne\\
&16&$181\ 16\ 26.5$&$20\ 51\ 19.3$&$-1.7$&$+0.5$&Dr.\ Maskelyne
\end{tabular}}
\end{center}

Aus dieser Vergleichung sieht man, dass man nach dieser ganzen Reihe von Beobachtungen nur wenig mehr an den Elementen würde ändern können, besonders da die letzten Greenwicher Declinationen im Mai, (welche einen grössern Einfluss auf die Elemente haben, als die geraden Aufsteigungen), noch sehr gut unter einander stimmen. Nur die berechneten Rectascensionen\HistoricalNote{20} müssten gegen das Ende hin noch etwa $4''$ kleiner gemacht werden; die Elemente würden aber dadurch nur ganz unerheblich afficirt werden. In einem Schreiben vom 26. Junius benachrichtigt uns Dr.\ \textsc{Gauss}, dass er die \textsc{Olbers}'schen Beobachtungen dieses Planeten vom 19.\ 20.\ und 21.\ Junius mit seinen III. Elementen auch noch verglichen habe, und diese stimmten noch so schön, dass er mit Gewissheit gar nichts darnach ändern könne. Bei diesen III. Elementen der Bahn wird man es demnach bewenden lassen können, bis die Beobachtungen ganz geschlossen sein werden. Damit berechnete demnach Dr.\ \textsc{Gauss} folgende
\Needspace{15\baselineskip}
\begin{center}
\textit{Ephemeride für die Pallas Olbersiana für Seeberger Mitternacht.}\par\medskip
{\small\setlength{\tabcolsep}{3pt}
\begin{tabular}{r r|r|r|r|r}
\multicolumn{2}{c|}{1802}&\shortstack{AR.\\der Pallas}&\shortstack{Declin.\\der Pallas}&\shortstack{Abstand\\von d.\ $\EarthSymbol$}&\shortstack{Abstand\\von d.\ $\odot$}\\\hline
Mai&18&$181^0\ 25'$&$20^0\ 55'$&$1.888$&$2.477$\\
&21&$181\ 40$&$20\ 59$&$1.930$&$2.484$\\
&24&$181\ 58$&$21\quad0$&$1.973$&$2.492$\\
&27&$182\ 19$&$20\ 59$&$2.016$&$2.500$\\
&30&$182\ 42$&$20\ 56$&$2.060$&$2.508$\\
Juni&2&$183\quad8$&$20\ 51$&$2.104$&$2.516$\\
&5&$183\ 36$&$20\ 44$&$2.148$&$2.524$\\
&8&$184\quad6\frac{1}{2}$&$20\ 36$&$2.192$&$2.532$
\end{tabular}
\par\medskip
\begin{tabular}{r r|r|r|r|r}
\multicolumn{2}{c|}{1802}&\shortstack{AR.\\der Pallas}&\shortstack{Declin.\\der Pallas}&\shortstack{Abstand\\von d.\ $\EarthSymbol$}&\shortstack{Abstand\\von d.\ $\odot$}\\\hline
Juni&11&$184^0\ 39'$&$20^0\ 26\frac{1}{2}'$&$2.237$&$2{,}539$\\
&14&$185\ 14$&$20\ 15$&$2.281$&$2.547$\\
&17&$185\ 51$&$20\quad3$&$2.326$&$2.555$\\
&20&$186\ 30$&$19\ 50$&$2.370$&$2.563$\\
&23&$187\ 10$&$19\ 35$&$2.415$&$2.571$\\
&26&$187\ 52\frac{1}{2}$&$19\ 19\frac{1}{2}$&$2.459$&$2.579$\\
&29&$188\ 36$&$19\quad3$&$2.504$&$2.587$
\end{tabular}}
\end{center}

Für Anfang des künftigen Jahres berechnete Dr.\ \textsc{Gauss} den Stand der \textit{Pallas} beiläufig also: den 31. Dec.\ um $12^{\mathrm{u}}$ Mitternacht in Seeberg Länge $253^0\ 33'$, Breite $26^0\ 10'$, nördl.\ Abstand von der Erde $3.779$, von der Sonne $3.043$. Allein es bleibt sehr ungewiss, ob wir diesen Planeten wegen seiner Lichtschwäche
\SourceLeaf{225}%
in den nächsten beiden Jahren wieder sehen werden. Dr.\ \textsc{Olbers} hat die beiden folgenden Oppositionen der \textit{Pallas} nach den \textsc{Gauss}'schen II Elementen berechnet. Als er sie den 28. März entdeckte, war ihr Abstand von der Erde nur $=1.38$; mit Ende Mai war er $=2.00$, und am 15. Juni $=2.31$. Künftiges Jahr 1803 gegen Ende des Juni kommt \textit{Pallas} wieder in die Opposition mit der Sonne; sie wird über dem Cerberus in $9^{\mathrm{z}}\ 8^0$ Länge und $48^0$ nördl.\ Breite zu stehen kommen, aber ihr Abstand von der Erde wird $2.61$ sein. Dr.\ \textsc{Olbers} glaubt, dass sie schwerlich in diesem Abstande bei der hellen nördlichen Dämmerung aufzufinden sein dürfte. Etwas besser geht es in der Opposition von 1804, die sich gegen das Ende des August in $11^{\mathrm{z}}\ 3^0$ mit $19^0$ N.\ Breite, also im Kopfe des Pegasus ereignet. Zwar findet Dr.\ \textsc{Olbers} den Abstand von der Erde noch $2.40$, aber die Nächte sind doch schon vollkommen dunkel. Aber \textit{ganz gewiss}, glaubt Dr.\ \textsc{Olbers}, werden wir die \textit{Pallas} im Jahr 1805 wieder sehen, da sie bei der Opposition schon in den der Sonne näher liegenden Theil ihrer Bahn herabgekommen ist, und also auch der Erde viel näher sein wird. Indessen werden hoffentlich diese Betrachtungen keinen Astronomen abhalten, die \textit{Pallas}, sobald die Möglichkeit der Aufsuchung vorhanden sein wird, nach Kräften aufzuspühren. Denn grosse und lichtstarke Fernröhre könnten denn doch wol hinreichend sein, diesen Planeten, und wenn er auch nur als ein Stern zwölfter Grösse erscheinen sollte, noch zu entdecken; freilich wird dessen Aufsuchung mit etwas mehr Schwierigkeit verbunden sein. Gesetzt aber auch, dass 1803 und 1804 alle Versuche fehlschlagen sollten, so hofft Dr.\ \textsc{Gauss} doch, dass seine Elemente der Bahn für das Jahr 1805 einen noch hinreichenden Leitfaden abgeben werden, um die \textit{Pallas}, (freilich mit etwas mehr Mühe als bei der \textit{Ceres}) dennoch wieder zu finden.\quad —\quad —\quad —

Was die Kleinheit dieser beiden neuen planetarischen Weltkörper betrifft, so ist es in der That schwer zu begreifen, wie ein so kleiner scheinbarer Durchmesser mit Genauigkeit gemessen werden kann. \textsc{Herschel} und \textsc{Schröter}, beide geübte Beobachter, und mit den grössten und besten Sehwerkzeugen versehen, haben solche Messungen versucht; dieser hat \textit{Ceres} 18 mal, \textit{Pallas} 40 mal grösser als jener gefunden. Um die Schwierigkeit solcher Messungen darzustellen, hat Dr.\ \textsc{Gauss} nach seinen letzten Elementen der Bahnen folgende Berechnungen der scheinbaren Durchmesser dieser Weltkörper angestellt.
\begin{center}
\textit{Durchmesser der Ceres.}\par\medskip
{\small\setlength{\tabcolsep}{6pt}
\begin{tabular}{r r|r|c}
\multicolumn{2}{c|}{1802}&\shortstack{Abstand\\von der\\$\EarthSymbol$}&\shortstack{scheinbarer Durchmesser, den wahren\\nach Dr.\ \textsc{Herschel} $=35$ Deutsche\\Meilen gesetzt.}\\\hline
März&28&$1.621$&$0''22$
\end{tabular}}
\par\bigskip
\textit{Durchmesser der Pallas.}\par\medskip
{\small\setlength{\tabcolsep}{6pt}
\begin{tabular}{r r|r|c}
\multicolumn{2}{c|}{1802}&\shortstack{Abstand\\von der\\$\EarthSymbol$}&\shortstack{Scheinb.\ Durchmesser, wenn man Dr.\\\textsc{Herschel}'s Resultat $=15$ Deutsche\\Meilen als richtig annimmt.}\\\hline
April&4&$1.407$&$0''11$\\
—&15&$1.386$&$0.10$\\
Juni&23&$2.415$&$0.06$
\end{tabular}}
\end{center}

Ob man solche kleine Durchmesser noch wirklich messen, selbst das reflectirte Licht von so kleinen Oberflächen noch sehen könne, überlassen wir andern zur Beurtheilung.\quad —\quad —\quad $=$
\begin{flushright}[\textsc{von Zach}.]\end{flushright}
\SourceLeaf{226}%
\bigskip\Needspace{9\baselineskip}
\pdfbookmark[1]{Pallas Olbersiana. 1802 August}{pallas180208}
\begin{center}
PALLAS OLBERSIANA.\par
\medskip\hrule\medskip
Monatliche Correspondenz. Band VI. S.\ 187—196. 1802 August.\par
\medskip\hrule
\end{center}\medskip
—\quad —\quad —\quad Auch aus \textit{Vilna}, \textit{Prag}, \textit{Cremsmünster} und \textit{Padua} haben wir die Bestätigung erhalten, dass die \textit{Pallas} an diesen Orten nicht beobachtet worden sei. Einer Nachricht des Dr.\ \textsc{Gauss} zu Folge hatte Dr.\ \textsc{Maskelyne} seine Beobachtung dieses Planeten am 20. Juli noch nicht aufgegeben, und glaubt die \textit{Pallas} den 18. Juli als einen Stern 11. Grösse noch beobachtet zu haben, wovon er jedoch nicht ganz sicher war, ob es der Planet oder ein Fixstern war, welches die Berechnung oder die folgende Beobachtung ausweisen wird.

Die III. Elemente der \textit{Pallas}-Bahn des Dr.\ \textsc{Gauss} stimmen noch fort\-während gut mit den beobachteten Declinationen; nur die geraden Aufsteigungen entfernen sich jetzt um einige Secunden mehr. Dr.\ \textsc{Gauss} verglich die in Mailand von \textsc{Cesaris} an einem vortrefflichen \textsc{Ramsden}'schen achtfüssigen Mauerquadranten angestellten Beobachtungen, welche wir im vorigen Hefte S.\ 73 mitgetheilt haben, und erhielt folgende Übereinstimmung:
\begin{center}
\textit{Vergleichung der Mailänder Meridian-Beobachtungen der Pallas}\par
\textit{mit den }\textsc{Gauss}\textit{'schen III. Elementen.}\par\medskip
{\small\setlength{\tabcolsep}{4pt}
\begin{tabular}{r r|r|r|r|r|r}
\multicolumn{2}{c|}{1802}&\shortstack{Mittl.\ Zeit\\in Mailand}&\multicolumn{2}{c|}{Berechnete}&\multicolumn{2}{c}{\shortstack{Unterschied\\R—B}}\\
\multicolumn{2}{c|}{}&&\shortstack{Gerade\\Aufsteigung}&\shortstack{Nördliche\\Abweichung}&in AR.&in Decl.\\\hline
Mai&4&$9^{\mathrm{u}}\ 15^{\mathrm{m}}\ 43^{\mathrm{s}}1$&$180^0\ 56'\ 15''7$&$20^0\quad2'\ 44''1$&$+15''7$&$-5''9$\\
&5&$9\ 11\ 45.3$&$180\ 55\ 56.8$&$20\quad8\ 38.1$&$+19.8$&$+5.1$\\
&6&$9\quad7\ 49.3$&$180\ 56\quad0.3$&$20\ 14\ 10.4$&$+16.3$&$+6.4$\\
&7&$9\quad3\ 55.6$&$180\ 56\ 24.8$&$20\ 19\ 22.1$&$+10.8$&$+2.1$\\
&8&$9\quad0\quad3.2$&$180\ 57\ 13.3$&$20\ 24\ 12.4$&$+8.3$&$+0.4$\\
&9&$8\ 56\ 12.1$&$180\ 58\ 21.9$&$20\ 28\ 43.1$&$+6.9$&$+0.1$\\
&10&$8\ 52\ 22.4$&$180\ 59\ 52.8$&$20\ 32\ 53.0$&$+7.8$&$-2.0$\\
&11&$8\ 48\ 34.3$&$181\quad1\ 45.2$&$20\ 36\ 43.4$&$+9.2$&$-3.6$\\
&12&$8\ 44\ 47.0$&$181\quad3\ 58.8$&$20\ 40\ 14{,}6$&$+7.8$&$-1.4$\\
&17&$8\ 26\ 12.9$&$181\ 20\ 17.1$&$20\ 53\ 18.0$&$+11.1$&$-4.0$\\
&18&$8\ 22\ 33.1$&$181\ 24\ 33.5$&$20\ 55\quad3.6$&$+15.5$&$+0.6$\\
&19&$8\ 18\ 55.8$&$181\ 29\quad9.4$&$20\ 56\ 30.8$&$+16.4$&$-0.2$\\
&20&$8\ 15\ 19.6$&$181\ 34\quad4.8$&$20\ 57\ 43.0$&$+14.8$&$-2.0$\\
&21&$8\ 11\ 44.6$&$181\ 39\ 18.6$&$20\ 58\ 39.5$&$+10.6$&$-1.5$\\
&22&$8\quad8\ 11.1$&$181\ 44\ 52.9$&$20\ 59\ 20.6$&$+6.9$&$-1.4$
\end{tabular}}
\end{center}

Auch Dr.\ \textsc{Olbers} Beobachtungen geben dieselbe Übereinstimmung. Die Harmonie derselben unter sich gibt zugleich den Beweis, welchen vorzüglichen Werth man diesen, obgleich nur am \textit{Kreismikrometer} angestellten Beobachtungen beilegen darf. Wer sich die Mühe nehmen will, sie so genau zu untersuchen, wie wir gethan haben, wird finden, dass sie den an Aequatorial-Sectoren angestellten nichts nachgeben, ihnen zur Seite stehen, mitunter auch vorzuziehen sind. Dieses beweist freilich mehr für die Geschicklichkeit des vortrefflichen Beobachters, als für die Güte des Werkzeuges, wie unsere Leser aus folgender Darstellung erkennen werden.
\newcommand{\PallasSymbol}{\text{\mdseries\AstroSymbols ⚴}}
\SourceLeaf{227}%
\Needspace{15\baselineskip}
\begin{center}
\textit{Vergleichung der Dr.\ }\textsc{Olbers}\textit{'schen Beobachtungen der Pallas}\par
\textit{mit den III.\ }\textsc{Gauss}\textit{'schen Elementen.}\par\medskip
{\small\setlength{\tabcolsep}{4pt}
\begin{tabular}{r r|r|r|r|r|r}
\multicolumn{2}{c|}{1802}&\shortstack{Mittl.\ Zeit\\in Bremen}&\multicolumn{2}{c|}{Berechnete}&\multicolumn{2}{c}{\shortstack{Unterschied\\R—B}}\\
\multicolumn{2}{c|}{}&&\shortstack{Gerade Auf-\\steig.\ der $\PallasSymbol$}&\shortstack{Gerade Auf-\\steig.\ der $\PallasSymbol$}&in AR.&in Decl.\\\hline
Juni&19&$11^{\mathrm{u}}\ 13^{\mathrm{m}}\ 55^{\mathrm{s}}$&$186^0\ 16'\ 19''$&$19^0\ 54'\ 17''$&$-6''$&$+17''$\\
&20&$11\quad7\ 51$&$186\ 29\ 12$&$19\ 49\ 39$&$-22$&$-1$\\
&21&$11\quad8\ 54$&$186\ 42\ 39$&$19\ 44\ 55$&$-20$&$-3$\\
&26&$11\ 18\ 34$&$187\ 52\quad3$&$19\ 19\ 31$&$-16$&$+8$\\
Juli&4&$11\ 28\ 14$&$189\ 52\ 12$&$18\ 33\ 39$&$-26$&$-5$\\
&8&$11\quad5\ 16$&$190\ 55\ 44$&$18\quad8\ 50$&$-25$&$-8$\\
&8&$11\ 13\ 20$&$190\ 55\ 50$&$18\quad8\ 48$&$-23$&$-10$\\
&9&$11\ 53\ 47$&$191\ 12\ 36$&$18\quad2\ 35$&$-28$&$+1$
\end{tabular}}
\end{center}

Aus allen diesen Vergleichungen und vortrefflichen Übereinstimmungen der Beobachtungen mit der \textsc{Gauss}'schen Ellipse (welche dieser geschickte Rechner der Wahrheit noch näher bringen wird) können wir der zuversichtlichen Hoffnung leben, dass wir mit dem Besitz dieses neuen Planeten selbst dann geborgen sein werden, wenn 1803 und 1804 seine Auffindung wegen seiner Lichtschwäche unmöglich werden sollte. Doch schöpfen wir schon gegenwärtig die gegründete Hoffnung, dass \textit{Pallas} künftiges Jahr sich unsern geschärften Blicken nicht entziehen wird. Dr.\ \textsc{Olbers}, welcher diesen Planeten noch zuletzt am 9. Juli gesehen hatte, ist selbst dieser Meinung und schreibt: „So schwer der Planet auch jetzt zu sehen ist, so bin ich doch überzeugt, dass Sie ihn künftiges Jahr an Ihrem Passagen-Instrumente sehen werden. Er ist alsdann in seiner Opposition der Erde etwas näher, als er ihr am 9. Juli war. Zwar steht er weiter von der Sonne und wird von dieser schwächer erleuchtet, allein dies wird überflüssig durch seine grössere Höhe über dem Horizont und die grössere Entfernung von der Dämmerung ersetzt.“

Auch Dr.\ \textsc{Gauss} hat die beste Hoffnung zur Wiederauffindung der \textit{Pallas} im künftigen Jahre. Er findet ihren Abstand von der Erde am 28. Jan.\ 1803 (etwa einen Tag vor der Opposition) $2.56$. Den Durchmesser dieses Planeten findet er aus Dr.\ \textsc{Herschel}'s eigener Messung etwas grösser. Den 22. April gibt Dr.\ \textsc{Herschel} nach einer ziemlich guten Messung den scheinbaren Durchmesser $=0''17$; hieraus berechnet Dr.\ \textsc{Gauss} den wahren Durchmesser $26\frac{1}{2}$ Deutsche Meile (den Abstand von der Erde $=1.562$). In seinem letzten Briefe äussert er sein Befremden über die so enorm verschiedenen Resultate von Dr.\ \textsc{Herschel}'s und Dr.\ \textsc{Schröter}'s Messungen dieser Durchmesser, die jedoch nach einerlei Methode sind gemacht worden. „Ich bin sehr neugierig, schreibt er, wie starke Vergrösserungen Dr.\ \textsc{Herschel} angewandt hat. 500malige Vergrösserung würde einen scheinbaren Durchmesser von $0''17$ wol noch nicht zur Scheibe machen?“ Ich wenigstens habe mit einer 300maligen Vergrösserung keine Spur von einer Scheibe weder an diesem \textsc{Olbers}'schen, noch an jenen \textsc{Piazzi}'schen Planeten wahrnehmen können. Auf eine Anfrage deshalb an Dr.\ \textsc{Olbers} erhielt ich die Antwort: „Auch ich habe nie mit den stärksten Vergrösserungen meines \textsc{Dollond}'schen Achromaten einen Unterschied zwischen diesen Planeten und Fixsternen von gleicher scheinbarer Grösse wahrnehmen können.“

Wir haben im vorigen Hefte S.\ 85 [S.\ 219 d.\ B.] eine \textsc{Gauss}'sche Ephemeride für die \textit{Pallas} mitgetheilt; hier folgt ihre Fortsetzung. Erscheint sie gleich zu spät, um noch die Beobachtungen erleichtern zu können: so wird sie doch dazu dienen, zu verhüten, dass etwa mit dem Planeten verwechselte Sterne nicht für die \textit{Pallas} ausgegeben werden. [Die Ephemeride ist ebenso eingerichtet, wie die auf S.\ 219 abgedruckte, enthält nur nicht den Abstand der Pallas von der Sonne und erstreckt sich bis 1802 Aug.\ 28. Die bis hieher in dieser Ausgabe von \textsc{Gauss} Werken mitgetheilten Ephemeriden werden als Beispiele der vom \textsc{Gauss} berechneten genügen, die übrigen aber zu Gunsten der Raumersparniss unberücksichtigt bleiben können.]
\expandafter\def\csname HistoricalNoteText21\endcsname{Die III. Elemente geben die AR.\ am 4. Febr.\ um $13'$ und am 28. Juni um $32'$ kleiner, die Decl.\ $2'$ grösser.}
\SourceLeaf{228}%
\bigskip\Needspace{10\baselineskip}
\pdfbookmark[1]{Gauss an Bode. 1802 Sept. 10}{bode18020910}
\begin{center}GAUSS AN BODE.\end{center}
\begin{flushright}Braunschweig 1802 Sept.\ 10.\end{flushright}
\begin{center}
\medskip\hrule\medskip
Astronomisches Jahrbuch für 1805. Berlin 1802. Seite 227.\par
\medskip\hrule
\end{center}\medskip

Die Elemente, nach denen die folgende Ephemeride für den 9ten Hauptplaneten: \textit{Pallas Olbersiana} berechnet ist, gründen sich auf die ersten Seeberger und letzten Greenwicher Merid.\ Beob.\ und Dr.\ \textsc{Olbers} Beob.\ vom 8. Juli und sind:
\begin{center}
{\small\setlength{\tabcolsep}{3pt}
\begin{tabular}{l r}
Epoche 1802 März 31 Seeberger Meridian $\ .\ .\ .\ .\ .$&$162^0\ 55'\quad6''8$\\
Tägliche tropische Bewegung $\ .\ .\ .\ .\ .\ .\ .\ .\ .$&$769''7263$\\
Logarithmus der halben grossen Axe $\ .\ .\ .\ .\ .\ .$&$0.4424992$\\
Sonnenferne $\ .\ .\ .\ .\ .\ .\ .\ .\ .\ .\ .\ .\ .\ .\ .$&$301^0\ 38'\ 42''$\\
Excentricität $\ .\ .\ .\ .\ .\ .\ .\ .\ .\ .\ .\ .\ .\ .\ .$&$0.243888$\\
Knoten {\AstroSymbols ☊} $\ .\ .\ .\ .\ .\ .\ .\ .\ .\ .\ .\ .\ .\ .\ .\ .$&$172^0\ 26'\ 31''$\\
Neigung $\ .\ .\ .\ .\ .\ .\ .\ .\ .\ .\ .\ .\ .\ .\ .\ .\ .$&$34\ 36\ 59$
\end{tabular}}
\end{center}
Wenn wir die Lichtstärke des Planeten am 8. Juli d.\ J.\ $=1$ setzen, so war sie
\begin{center}
\begin{tabular}{r l r}
am 4. April&$\ .\ .\ .\ .\ .\ .\ .\ .\ .\ .$&$=4.23$\\
4. Mai&$\ .\ .\ .\ .\ .\ .\ .\ .\ .\ .$&$=2.75$\\
26. Juli&$\ .\ .\ .\ .\ .\ .\ .\ .\ .\ .$&$=0.80$\\
24. August&$\ .\ .\ .\ .\ .\ .\ .\ .\ .$&$=0.60$\\
\multicolumn{2}{l}{grösste Lichtstärke im nächsten Jahre}&$=0.653$
\end{tabular}
\end{center}

Dies zeigt, dass die Hoffnung, den Planeten 1803 wieder zu finden, noch nicht ganz zuverlässig sein kann. Die grosse Höhe über dem Horizont und die grössere Entfernung von der Dämmerung müssen das Abgehende ersetzen. Doch schmeichle ich mir, dass in den südlichen Klimaten, z.\ B.\ in Palermo, wo die Juninächte um Mitternacht ganz dunkel sind, die Auffindung leichter gelingen werde.

Wenn ich künftig auch eine Ephemeride der $\CeresSymbol$ nach verbesserten Elementen berechnen werde, so werde ich nicht ermangeln, dieselbe Ew.\ — gleichfalls zuzusenden\HistoricalNote{21}. [Die Ephemeride geht von Febr.\ 4 bis Juni 28 Mitternacht im Seeberger Meridian, 1803, und gibt mit Intervallen von drei Tagen AR.\ und Decl.\ bis auf Minuten genau.]
\begin{center}
\begin{tabular}{l}
{\AstroSymbols ☍} Juni 30. Nachmittag.\\
nach den III. Elementen Juni 30. Vormittag bürgerl.\ Zeit.
\end{tabular}
\end{center}
\SourceLeaf{229}%
\bigskip\Needspace{9\baselineskip}
\pdfbookmark[1]{Ceres Ferdinandea. 1802 Juni}{ceres180206}
\begin{center}
CERES FERDINANDEA.\par
\medskip\hrule\medskip
Monatliche Correspondenz. Band V. S.\ 576—590. 1802 Juni.\par
\medskip\hrule
\end{center}\medskip
—\quad —\quad —

Auch Dr.\ \textsc{Gauss} hat, wie bisher, fortgefahren, unsere sämmtlichen Beobachtungen dieses Planeten mit den VII. Elementen seiner Bahn zu vergleichen. Diese Vergleichung bestätigt nicht allein, dass die Fehler dieser Elemente sich sehr wenig und langsam ändern, folglich den Ort des Planeten während der ganzen diesjährigen Sichtbarkeit überflüssig genau darstellen, sondern diese Vergleichung wird ihm auch künftig noch, bei einer feinern Correction dieser Elemente, und wenn die Störungs-Gleichungen mitgenommen werden, sehr zu Statten kommen. Bis zum 19. März hatte Dr.\ \textsc{Gauss} im Mai-Hefte S.\ 470 [S.\ 211 d.\ B.] diesen Vergleich fortgeführt; hier folgen nun die Resultate zwölf darauf folgender Beobachtungen.
\begin{center}
{\small\setlength{\tabcolsep}{4pt}
\begin{tabular}{r r|r|r|r|r|r}
\multicolumn{2}{c|}{1802}&\shortstack{Mittl.\ Zeit\\in Seeberg}&\shortstack{Berechnete AR.\\der Ceres}&\shortstack{Berechnete\\Abweichung\\der Ceres}&\multicolumn{2}{c}{\shortstack{Unterschied\\R—B}}\\
\multicolumn{2}{c|}{}&&&&in AR.&in Decl.\\\hline
März&23&$12^{\mathrm{u}}\quad6^{\mathrm{m}}\ 28^{\mathrm{s}}9$&$182^0\ 20'\ 39''9$&$17^0\ 36'\ 51''4$&$+16''1$&$+23''5$\\
&27&$11\ 47\ 20.3$&$181\ 29\ 21.2$&$17\ 50\ 55.5$&$+20.1$&$+25.9$\\
&28&$11\ 42\ 33.7$&$181\ 16\ 41.4$&$17\ 53\ 52.9$&$+23.7$&$+29.8$\\
&29&$11\ 37\ 48.0$&$181\quad4\quad7.5$&$17\ 56\ 36.3$&$+18.3$&$+32.3$\\
&30&$11\ 33\quad2.5$&$180\ 51\ 40.3$&$17\ 59\quad5.5$&$+16.4$&$+36.5$\\
&31&$11\ 28\ 17.0$&$180\ 39\ 21.1$&$18\quad1\ 20.3$&$+24.1$&$+31.9$\\
April&1&$11\ 23\ 32.4$&$180\ 27\ 10.7$&$18\quad3\ 20.6$&$+25.2$&$+30.8$\\
&2&$11\ 18\ 48.7$&$180\ 15\ 10.1$&$18\quad5\quad6.2$&$+23.6$&$+31.4$\\
&3&$11\ 14\quad6.0$&$180\quad3\ 20.1$&$18\quad6\ 37.4$&$+17.7$&$+27.0$\\
&4&$11\quad9\ 24.1$&$179\ 51\ 41.7$&$18\quad7\ 53.5$&$+10.8$&$+24.1$\\
&5&$11\quad4\ 41.9$&$179\ 40\ 15.6$&$18\quad8\ 54.9$&$+21.5$&$+27.6$\\
&7&$10\ 55\ 21.3$&$179\ 18\quad3.9$&$18\ 10\ 12.5$&$+24.2$&$+25.4$
\end{tabular}}
\end{center}

Den Gegenschein des Planeten mit der Sonne berechnete Dr.\ \textsc{Gauss} aus unseren Beobachtungen vom 15.\ 16.\ 17.\ 18.\ 19.\ März, und fand, dass er sich den 17. März um $4^{\mathrm{u}}\ 21^{\mathrm{m}}\ 10^{\mathrm{s}}$ mittl.\ Seeberger Zeit ereignet habe, in $5^{\mathrm{s}}\ 26^0\ 21'\ 27''3$ der Länge und $17^0\quad7'\ 59''6$ der geocentr.\ Breite.
\par\noindent —\quad —\quad —\hfill [\textsc{von Zach}.]\par

\bigskip\Needspace{9\baselineskip}
\pdfbookmark[1]{Ceres Ferdinandea. 1802 October}{ceres180210}
\begin{center}
CERES FERDINANDEA.\par
\medskip\hrule\medskip
Monatliche Correspondenz. Band VI. S.\ 382—389. 1802 October.\par
\medskip\hrule
\end{center}\medskip
—\quad —\quad —

Die Störungen der Länge und des Abstandes der \textit{Ceres} durch Jupiter sind zwar schon von mehreren Astronomen in Rechnung genommen worden: allein solche verwickelte und schwierige Rechnungen können nicht genug wiederholt werden. Erst kürzlich schrieb uns Prof.\ \textsc{Wurm} über diesen Gegenstand, und bei Gelegenheit einiger eingeschickten Verbesserungen und erläuternden Zusätze zu den Formeln der
\newcommand{\JupiterSymbol}{\text{\mdseries\AstroSymbols ♃}}
\SourceLeaf{230}%
Mars-Störung, welche wir nächstens in unsern Heften mittheilen werden: Ich glaube überhaupt, wem auf diesem Wege der Störungs-Rechnungen nie Dornen vorgekommen sind, der hat gewiss den Weg selbst nie betreten, denn dass man bei solchen Untersuchungen hin und wieder unter Dornen geräth, und selbst auch etwa, bis man sich besser orientirt hat, ein kleines \textit{quid pro quo} setzt, ist nicht wol zu vermeiden.

Dr.\ \textsc{Gauss} hat daher diese Störungs-Rechnungen nach seinen VII. Elementen der Ceres-Bahn wiederholt, und hier und da kleine Unterschiede gefunden. Die analytischen Formeln, die er bei dieser Rechnung gebraucht, hat er sich alle erst selbst entwickelt; sie sind zum Theil von den \textsc{La Place}schen in der Form etwas verschieden ausgefallen.
\begin{center}
\textit{Störungen der Ceres Ferdinandea durch Jupiter, von Dr.\ }\textsc{Gauss}\textit{ berechnet.}\par\medskip
{\small\setlength{\tabcolsep}{3pt}
\begin{tabular}{l r}
Jährliche Abnahme der Excentricität $\ .\ .\ .\ .\ .\ .\ .\ .\ .\ .\ .\ .$&$0.000005909$\\
Jährliche Bewegung der Sonnenferne gegen die Fixsterne $\ .\ .\ .\ .\ .$&$+70''15$
\end{tabular}}
\par\bigskip
\textit{Periodische Gleichungen.}\par\medskip
1)\quad \textbf{Für die Länge in der Bahn.}
\end{center}
\[
\begin{array}{l}
-231''94\sin(\CeresSymbol-\JupiterSymbol)\\
+496.68\sin2(\CeresSymbol-\JupiterSymbol)\\
+44.15\sin3(\CeresSymbol-\JupiterSymbol)\\
+10.07\sin4(\CeresSymbol-\JupiterSymbol)\\
+3.05\sin5(\CeresSymbol-\JupiterSymbol)\\
+1.07\sin6(\CeresSymbol-\JupiterSymbol)\\
+0.41\sin7(\CeresSymbol-\JupiterSymbol)\\
-60.25\sin(\JupiterSymbol+17^0\ 41'\ 28'')\\
-621.29\sin(\CeresSymbol-2\JupiterSymbol-26^0\ 49'\ 22'')
\end{array}
\]
\[
\begin{array}{l}
-442''65\sin(2\CeresSymbol-3\JupiterSymbol-11^0\ 34'\ 39'')\\
+56.69\sin(3\CeresSymbol-4\JupiterSymbol-11\quad24\quad1'')\\
+10.58\sin(4\CeresSymbol-5\JupiterSymbol-11\quad11\quad21'')\\
+3.26\sin(5\CeresSymbol-6\JupiterSymbol-10\quad55\quad43'')\\
+23.62\sin(2\CeresSymbol-\JupiterSymbol+36\quad31\quad34'')\\
-53.93\sin(3\CeresSymbol-2\JupiterSymbol+33\quad24\quad19'')\\
-5.96\sin(4\CeresSymbol-3\JupiterSymbol+31\quad20\quad18'')\\
-1.70\sin(5\CeresSymbol-4\JupiterSymbol+29\quad33\quad57'')
\end{array}
\]
\begin{center}2)\quad \textbf{Für den Radius vector.}\end{center}
\[
\begin{array}{l}
-0.0000947\\
+0.0010304\cos(\CeresSymbol-\JupiterSymbol)\\
-0.0038023\cos2(\CeresSymbol-\JupiterSymbol)\\
-0.0004206\cos3(\CeresSymbol-\JupiterSymbol)\\
-0.0001077\cos4(\CeresSymbol-\JupiterSymbol)\\
-0.0000350\cos5(\CeresSymbol-\JupiterSymbol)\\
-0.0000128\cos6(\CeresSymbol-\JupiterSymbol)\\
-0.0000050\cos7(\CeresSymbol-\JupiterSymbol)\\
-0.0000633\cos(\CeresSymbol+25^0\ 56'\ 46'')\\
+0.0002469\cos(\JupiterSymbol+23\quad42\quad53'')
\end{array}
\]
\[
\begin{array}{l}
+0.0008589\cos(\CeresSymbol-2\JupiterSymbol-24^0\ 37'\ 54'')\\
+0.0025827\cos(2\CeresSymbol-3\JupiterSymbol-11\quad19\quad13'')\\
-0.0004817\cos(3\CeresSymbol-4\JupiterSymbol-11\quad32\quad33'')\\
-0.0001075\cos(4\CeresSymbol-5\JupiterSymbol-11\quad31\quad15'')\\
-0.0000366\cos(5\CeresSymbol-6\JupiterSymbol-11\quad21\quad15'')\\
-0.0001367\cos(2\CeresSymbol-\JupiterSymbol+37\quad40\quad22'')\\
+0.0003049\cos(3\CeresSymbol-2\JupiterSymbol+32\quad54\quad59'')\\
+0.0000411\cos(4\CeresSymbol-3\JupiterSymbol+29\quad43\quad48'')\\
+0.0000139\cos(5\CeresSymbol-4\JupiterSymbol+27\quad40\quad24'')
\end{array}
\]
\begin{flushright}[\textsc{von Zach}.]\end{flushright}
\SourceLeaf{231}%
\bigskip\Needspace{9\baselineskip}
\pdfbookmark[1]{Pallas Olbersiana. 1802 October}{pallas180210}
\begin{center}
PALLAS OLBERSIANA.\par
\medskip\hrule\medskip
Monatliche Correspondenz. Band VI. S.\ 390—396. 1802 October.\par
\medskip\hrule
\end{center}\medskip
—\quad —\quad —

Dr.\ \textsc{Gauss}, welcher seine Elemente der Pallas-Bahn zum drittenmale verbessert hatte, hat nun eine \textit{vierte} Correction gewagt. Diese neuen Elemente sind aus unseren ersten \textit{Seeberger}, aus Dr.\ \textsc{Maskelyne}'s letzten \textit{Greenwicher} Meridian-Beobachtungen, und aus Dr.\ \textsc{Olbers} Beobachtung vom 8. Juli abgeleitet. Dr.\ \textsc{Gauss} will sie indessen noch eben nicht für sicherer als die III. Elemente ausgeben, da er eine viel genauere Verbesserung von den spätern Mailänder Beobachtungen erwartet; diese Elemente sind indessen folgende: [dieselben wie die oben in dem Briefe an \textsc{Bode} vom 10. Sept.\ S.\ 223 angegebenen, ebenfalls für die Epoche 1802 März 31 Mittag in Seeberg mit Hinzufügung bei Sonnenferne und Knoten ‘für die Epoche, siderisch ruhend’].

Mit diesen Elementen hat Dr.\ \textsc{Gauss} nachfolgende Ephemeride für diesen Planeten auf künftiges Jahr voraus berechnet. Es ist nützlich, sie noch bei guter Zeit bekannt zu machen, weil man indessen in der Gegend, welche \textit{Pallas} im J.\ 1803 durchwandern wird, zweifelhafte Sterne genauer bestimmen, sich in dieser Himmelsgegend im voraus orientiren und bekannt machen kann, wodurch uns das Auffinden mehr erleichtert werden wird. In der letzten Columne der Ephemeride, welche die Lichtstärke des Planeten enthält, ist diejenige zur Einheit angenommen worden, die der Planet am 8. Juli dieses Jahres hatte; nach eben diesem Massstabe war sie:
\begin{center}
\begin{tabular}{r l l r}
am 4. April&1802&$=$&$4.23$\\
4. Mai&—&&$2.75$\\
16. Juli&—&&$0.90$\\
26. Juli&—&&$0.80$\\
7. Aug.&—&&$0.70$\\
24. Aug.&—&&$0.60$
\end{tabular}
\end{center}
Hiebei ist aber die Grösse der Phase, die Höhe über dem Horizont und die Entfernung von der Dämmerung nicht in Betrachtung gezogen. Der erste Umstand ist ganz unerheblich, die beiden andern hängen von den Beobachtungsorten und den mächtigeren Fernröhren ab. Die grösste Lichtstärke, welche \textit{Pallas} im künftigen Jahre erreicht, wird $=0.656$ sein; dieselbe hatte sie in diesem 1802 Jahre am 14. Aug.\ \textsc{Oriani} \textit{beobachtete} diesen Planeten noch am 8. Aug., und \textit{sah} ihn mit Mühe den 17. und 18. dieses Monats, wozu aber wahrscheinlich die Dämmerung und der tiefe Stand des Planeten vieles beitrug. Es bleibt uns hienach noch immer die angenehme Hoffnung, dass wir künftiges Jahr die \textit{Pallas} zu Gesichte bekommen, und ihren Gegenschein mit der Sonne werden beobachten können, da wahrscheinlich der Planet schon vor dieser Epoche (30 Juni 1803) aufgefunden sein wird; denn mit der Zusammenkunft nimmt seine Sichtbarkeit oder Lichtstärke schon wieder ab. Im Meridian werden wir wol, wenigstens in unsern Breiten, vor der letzten Hälfte des Mai oder Anfang Juni wenig Hoffnung haben, diesen Planeten zu beobachten. [Die Ephemeride geht von Febr.\ 4. bis Juni 28. Mitternacht Seeberg 1803 gibt für jeden dritten Tag AR.\ und Abweichen bis auf Minuten genau, den Abstand von Erde in zwei, die Lichtstärke in drei Decimalstellen.]

Sollte es dem Dr.\ \textsc{Gauss} auch gelingen, wie es höchst wahrscheinlich ist, obige Elemente durch die Italienische Beobachtung noch zu verbessern: so wird es doch nicht nöthig sein, die beschwerliche
